% sec-06-cmc.tex - Chemistry, Manufacturing and Control (full stage)
\section{Chemistry, Manufacturing and Control Information}

\subsection{Chemistry, Manufacturing and Control}

This section presents the Chemistry, Manufacturing and Control (CMC)
information for the investigational drug daraxonrasib (RMC-6236) supporting the
Phase 1, first-in-human, open-label, single-arm investigation described
elsewhere in this Investigational New Drug application (IND v1.0, repository
v4.3.0). It is organized and submitted in accordance with 21 CFR \S
312.23(a)(7), which requires sufficient information about the composition,
manufacture, and control of the drug substance and the drug product to assure
the proper identification, quality, purity, and strength of the investigational
drug, scaled to the early phase of investigation, the limited number of
participants, and the short duration of exposure characteristic of a Phase 1
study \cite{phase1guidance}. Consistent with the FDA guidance on the content
and format of CMC information for a Phase 1 study, the emphasis here is on the
information needed to assure participant safety, including identity, strength,
quality, purity, and stability, rather than on the detailed validation data
appropriate to later phases \cite{phase1guidance}. As the investigation
proceeds and the manufacturing process and analytical methods mature, the CMC
information will be updated through information amendments under 21 CFR \S
312.31 as required.

The investigational regimen couples this oral drug with an on-premises large
language model (LLM) directed eight-arm robotic pancreaticoduodenectomy
(Whipple) platform that is the subject of the parallel investigational device
exemption content of this submission. The device is a durable instrument and is
not a drug, a biologic, or a consumable manufactured lot; its design,
specification, calibration, and quality controls are addressed in the device
description and the Physical AI System Description fields of 21 CFR \S
312.23(g), not in this CMC section
\cite{daraxwhipple,onpremwhippl}. Accordingly, the four CMC subsections below
address only the orally administered drug substance and drug product, the
non-applicable placebo, and labeling, followed by the environmental assessment.
The CMC information architecture and the mapping of each subsection to its
content is shown in Figure 16.

\begin{mermaidfig}
\node[mmgoal] (cmc) at (5.2,0)      {Chemistry, Manufacturing\\and Control ($\S$7)};
\node[mmproc] (ds)  at (0,-2.6)     {7.1.1 Drug substance\\daraxonrasib\\(RMC-6236)};
\node[mmproc] (dp)  at (3.6,-2.6)   {7.1.2 Drug product\\oral tablet,\\160 / 220 / 300 mg};
\node[mmctx]  (pl)  at (7.0,-2.6)   {7.1.3 Placebo product\\not applicable\\(open-label)};
\node[mmproc] (lb)  at (10.4,-2.6)  {7.1.4 Labeling\\``Caution:\\investigational''\\21 CFR 312.6};
\node[mmproc] (ea)  at (12.6,-0.2)  {7.2 Environmental\\assessment /\\categorical\\exclusion 25.31};
\node[mmin]   (ds1) at (0,-5.4)     {Characterization,\\specification,\\impurities};
\node[mmin]   (dp1) at (3.6,-5.4)   {Composition,\\container-closure,\\child-resistant};
\node[mmin]   (st)  at (7.0,-5.4)   {Stability program\\(supports shelf life)};
\draw[mmedge] (cmc) -- (ds);
\draw[mmedge] (cmc) -- (dp);
\draw[mmedge] (cmc) -- (pl);
\draw[mmedge] (cmc) -- (lb);
\draw[mmedge] (cmc.east) to[out=0,in=90,looseness=0.6] (ea.north);
\draw[mmedge] (ds) -- (ds1);
\draw[mmedge] (dp) -- (dp1);
\draw[mmedge] (dp) to[out=-90,in=90,looseness=0.7] (st.north);
\end{mermaidfig}
\figcaption{Figure 16. The Chemistry, Manufacturing and Control information
architecture for $\S$7. The drug substance (daraxonrasib, RMC-6236) is described
by characterization, specification, and impurity controls; the drug product is
the oral tablet at the 160, 220, and 300 mg strengths with its composition and
child-resistant, light-protective container-closure under 21 CFR $\S$312.6,
supported by a stability program; the placebo product is not applicable for this
open-label study; labeling carries the investigational caution statement; and an
environmental assessment or categorical exclusion under 21 CFR $\S$25.31 is
provided.}

\subsubsection{Drug Substance}

\paragraph{Nomenclature, description, and pharmacological class.}
The drug substance is daraxonrasib, also designated by the developmental code
RMC-6236, an orally bioavailable small molecule of the RAS(ON)
multi-selective (pan-RAS) inhibitor class. Daraxonrasib acts by binding the
active, guanosine triphosphate (GTP) bound state of mutant RAS in a tri-complex
with the endogenous chaperone cyclophilin A, thereby occluding the
RAS to effector interface and suppressing downstream mitogen-activated protein
kinase and phosphoinositide 3-kinase signaling across the common KRAS
oncoprotein variants found in pancreatic ductal adenocarcinoma (PDAC), including
the G12D, G12V, and G12R substitutions targeted by this protocol
\cite{daraxwhipple,pdac060s2030}. The agent received FDA Breakthrough Therapy
designation in June 2025 for previously treated metastatic PDAC bearing a KRAS
G12 mutation, and 300 mg administered orally once daily is the recommended
Phase 2 dose established in the pan-KRAS RASolute 302 program to which the dose
and exposure assumptions of this protocol are anchored
\cite{daraxwhipple,pdac060s2030}. The drug substance is a non-biologic, chirally
defined organic small molecule; it is not a radioactive drug, is not a
controlled substance, and is not a peptide, protein, or nucleic acid. The drug
substance is presented as a solid, and in the finished product it is formulated
as an immediate-release oral tablet.

\paragraph{Manufacturer and method of preparation.}
The drug substance is manufactured by, or under contract to, the sponsor's
designated active pharmaceutical ingredient manufacturer under current good
manufacturing practice (cGMP) appropriate to the investigational stage, and is
released against the drug-substance specification summarized in Table 7. The
synthetic route is a defined, multi-step convergent organic synthesis with
isolated, characterized intermediates and controlled reagents, catalysts, and
solvents; in-process controls govern the critical steps, and the final isolation
and drying steps establish the solid-state form, particle size distribution, and
residual-solvent profile of the substance. The name, address, and
establishment information of the drug-substance manufacturer, the detailed
reaction scheme, the list of starting materials and reagents, the in-process
controls, and the batch records are provided in the supporting CMC documentation
on file with the sponsor and are cross-referenced here consistent with the
Phase 1 level of CMC detail \cite{phase1guidance}. A statement that the drug
substance is manufactured in compliance with cGMP, and the certificate of
analysis for the clinical batch, accompany each shipment.

\paragraph{Characterization.}
The structure and identity of daraxonrasib are established by a panel of
orthogonal physicochemical methods. Structural elucidation and confirmation use
proton and carbon nuclear magnetic resonance spectroscopy, high-resolution mass
spectrometry, infrared spectroscopy, and ultraviolet to visible spectroscopy;
elemental analysis confirms empirical composition. Chiral purity is confirmed by
chiral high-performance liquid chromatography (HPLC) to assure the correct
stereochemical configuration of the defined stereocenters. The solid-state form
(polymorph) is characterized by X-ray powder diffraction and differential
scanning calorimetry, and the manufacturing process is controlled to deliver the
intended crystalline form, because solid-state form governs dissolution and
hence the oral exposure on which the pharmacokinetically gated restart advisory
of the protocol depends. Hygroscopicity, aqueous and physiological-pH
solubility, the partition coefficient, and the pKa are determined to support the
biopharmaceutic understanding of the tablet. These characterization data
constitute the reference standard against which clinical batches are compared.

\paragraph{Specification and control of impurities.}
The drug-substance specification, with tests, analytical procedures, and
acceptance criteria appropriate to a first-in-human Phase 1 study, is given in
Table 7. Organic impurities (process-related and degradation-related), residual
solvents, and elemental impurities are controlled to limits consistent with the
applicable International Council for Harmonisation (ICH) quality guidelines (ICH
Q3A for impurities in new drug substances, ICH Q3C for residual solvents, and
ICH Q3D for elemental impurities), as adapted to the limited duration and
participant exposure of a Phase 1 investigation
\cite{ichpaistduni,phase1guidance}. Any individual or total impurity exceeding
its qualification threshold is qualified by reference to the nonclinical
toxicology program summarized in the pharmacology and toxicology section of this
IND. Genotoxic-impurity risk is assessed and controlled under the ICH M7
framework, with structural alerts evaluated and relevant potential mutagenic
impurities controlled to the threshold of toxicological concern. The acceptance
criteria below are interim Phase 1 specifications and will be tightened as
process knowledge and additional batch data accumulate.

\begin{table}[t]
\centering
\caption{Drug-substance specification for daraxonrasib (RMC-6236): interim
Phase 1 release and stability tests, analytical procedures, and acceptance
criteria.}
\label{tab:cmc-ds-spec}
\begin{tabularx}{\textwidth}{L{4.1cm} L{4.0cm} Y}
\toprule
\textbf{Attribute / test} & \textbf{Analytical procedure} & \textbf{Acceptance criterion (interim Phase 1)} \\
\midrule
Appearance & Visual & White to off-white solid; free of foreign matter \\
Identity (spectroscopic) & IR; UV to visible & Conforms to reference standard \\
Identity (chromatographic) & HPLC retention time & Retention time matches reference standard \\
Identity (structural) & 1H and 13C NMR; HRMS & Consistent with assigned structure \\
Stereochemical purity & Chiral HPLC & Single defined stereoisomer; undesired enantiomer not more than 0.5 percent \\
Assay (content) & HPLC with UV detection & 98.0 to 102.0 percent on the anhydrous, solvent-free basis \\
Organic impurities (each) & HPLC, related substances & Not more than 0.5 percent per specified or unspecified impurity (ICH Q3A) \\
Organic impurities (total) & HPLC, related substances & Not more than 2.0 percent total \\
Genotoxic impurities & LC-MS or GC-MS, ICH M7 & Below the threshold of toxicological concern for specified mutagenic impurities \\
Residual solvents & Headspace GC & Within ICH Q3C Class limits for solvents used in synthesis \\
Elemental impurities & ICP-MS & Within ICH Q3D oral permitted daily exposure limits \\
Water content & Karl Fischer titration & Not more than 1.0 percent w/w \\
Residue on ignition / sulfated ash & USP general chapter & Not more than 0.2 percent \\
Solid-state form & X-ray powder diffraction & Conforms to the designated reference polymorph \\
Particle size distribution & Laser diffraction & Within the validated process range supporting dissolution \\
Microbial limits & USP total aerobic / yeast and mold & Within USP limits for an oral non-sterile substance \\
\bottomrule
\end{tabularx}
\figcaption{Table 7. Interim Phase 1 drug-substance specification. Acceptance
criteria are scaled to the early phase of investigation and will be tightened as
batch and process data accumulate; impurity limits reference ICH Q3A, Q3C, Q3D,
and M7 as adapted for Phase 1.}
\end{table}

\subsubsection{Drug Product}

\paragraph{Description and composition.}
The drug product is an immediate-release, film-coated oral tablet supplied in the
three strengths required to compose the once-daily doses of the 3+3 escalation:
160 mg for dose level 1 (DL1), 220 mg for dose level 2 (DL2), and 300 mg for
dose level 3 (DL3), the last of which corresponds to the RASolute 302 recommended
Phase 2 dose \cite{daraxwhipple,pdac060s2030}. The three strengths are
dose-proportional, sharing a common qualitative composition and a common blend so
that the only quantitative difference among strengths is the content of drug
substance and the corresponding fill weight, with the tablets differentiated by
size, debossing, or both to prevent dose confusion in this open-label study. The
qualitative and representative quantitative composition is given in Table 8. The
excipients are standard, compendial pharmaceutical excipients used at
conventional levels; the formulation contains a filler or diluent, a binder, a
disintegrant, a glidant, a lubricant, and a non-functional aqueous film coat. No
excipient of human or animal origin presenting a transmissible spongiform
encephalopathy risk is used, and no novel excipient is included.

\begin{table}[t]
\centering
\caption{Representative qualitative and quantitative composition of the
daraxonrasib immediate-release oral tablet (per the 160 mg strength; the 220 mg
and 300 mg strengths are dose-proportional from a common blend).}
\label{tab:cmc-dp-comp}
\begin{tabularx}{\textwidth}{L{4.4cm} L{4.6cm} Y}
\toprule
\textbf{Component} & \textbf{Function} & \textbf{Quantity (160 mg strength) and notes} \\
\midrule
Daraxonrasib (RMC-6236) & Active pharmaceutical ingredient & 160 mg per tablet (220 mg / 300 mg in higher strengths) \\
Microcrystalline cellulose & Filler / diluent & Quantity sufficient to target fill weight \\
Lactose monohydrate or mannitol & Filler / diluent & Compendial grade, conventional level \\
Povidone or hypromellose & Binder & Conventional level \\
Croscarmellose sodium & Disintegrant & Conventional level for immediate release \\
Colloidal silicon dioxide & Glidant & Conventional level \\
Magnesium stearate & Lubricant & Conventional level (typically not more than 1.0 percent) \\
Film coat (hypromellose-based) & Aqueous non-functional coat & Applied to target weight gain; removed as water on drying \\
Purified water & Coating vehicle & Removed during processing (not present in finished tablet) \\
\bottomrule
\end{tabularx}
\figcaption{Table 8. Drug-product composition. All excipients are compendial and
used at conventional levels; the film coat is non-functional and cosmetic. The
220 mg and 300 mg strengths are manufactured dose-proportionally from a common
blend, differing only in drug-substance content and fill weight and in tablet
size or debossing for identification.}
\end{table}

\paragraph{Manufacture.}
The drug product is manufactured under cGMP appropriate to the investigational
stage by the sponsor's designated drug-product manufacturer. The process is a
conventional solid-oral-dosage process: dispensing and screening of the drug
substance and excipients, blending (with wet or dry granulation as defined in the
batch record), lubrication, compression of the final blend into core tablets on a
rotary tablet press under in-process controls for tablet weight, hardness,
thickness, and friability, and aqueous film coating of the cores in a perforated
coating pan to a target weight gain, followed by drying. In-process controls at
compression include tablet weight and weight uniformity, hardness, thickness,
disintegration, and visual inspection; the coating step is controlled for weight
gain and appearance. The batch formula, the master and executed batch records,
the in-process control limits, and the name and address of the drug-product
manufacturer are provided in the supporting CMC documentation and are
cross-referenced here at the Phase 1 level of detail \cite{phase1guidance}. A
certificate of analysis accompanies each clinical batch, and each batch is
released against the drug-product specification before use.

\paragraph{Drug-product specification.}
Each clinical batch is tested and released against a drug-product specification
that includes appearance, identity of the drug substance (by HPLC retention time
and a spectroscopic method), assay (95.0 to 105.0 percent of label claim by HPLC),
content uniformity by the USP uniformity-of-dosage-units test, degradation
products controlled to interim Phase 1 limits consistent with ICH Q3B (not more
than 0.5 percent per specified or unspecified degradant and not more than 2.0
percent total), dissolution by the USP apparatus and medium validated for the
product (for example, not less than 80 percent dissolved in 30 minutes as a
provisional immediate-release criterion to be confirmed by the developing
dissolution method), water content, and microbial limits within USP acceptance
criteria for an oral non-sterile product. The dissolution method and acceptance
criterion are provisional and will be finalized as biopharmaceutic and stability
data accrue.

\paragraph{Container-closure system.}
The drug product is packaged in a child-resistant, light-protective,
moisture-protective container-closure system, namely high-density polyethylene
bottles with a child-resistant, induction-sealed closure and an appropriate
desiccant, or equivalent child-resistant unit-dose blister packaging. The
container-closure components contacting the product are of compendial
pharmaceutical grade and are suitable for an oral solid dosage form; the
child-resistant feature satisfies the special-packaging expectation for an
investigational oral oncology agent, and the light-protective and
moisture-protective features support the stability claim. Packaging and labeling
operations are performed under cGMP, and the investigational caution statement and
the label content described in Table 9 and in $\S$7.1.4 appear on the immediate
and any secondary container in accordance with 21 CFR \S 312.6 \cite{phase1guidance}.

\paragraph{Stability program.}
Stability of the drug product is supported by an ongoing stability program
conducted under ICH Q1A conditions on representative batches in the proposed
container-closure system, using stability-indicating analytical methods. The
program comprises long-term, intermediate, and accelerated conditions, with
attributes monitored at each station including appearance, assay, degradation
products, dissolution, and water content; a photostability evaluation under ICH
Q1B and an in-use stability evaluation for the multi-dose bottle presentation are
included. The program and its current status are summarized in Table 10. The
provisional shelf life (expiry or retest period) assigned to the clinical drug
product is supported by the available long-term and accelerated data and is
applied through expiry dating on the label; product is dispensed only within the
labeled expiry or retest date. Stability is monitored throughout the
investigation, and the shelf life will be extended through information amendments
as additional long-term data become available \cite{phase1guidance}. The drug
substance is also placed on a parallel ICH Q1A stability program to establish its
retest period.

\begin{table}[t]
\centering
\caption{Drug-product stability program for the daraxonrasib oral tablet:
storage conditions, time points, attributes monitored, and purpose.}
\label{tab:cmc-stability}
\begin{tabularx}{\textwidth}{L{3.5cm} L{3.3cm} L{4.2cm} Y}
\toprule
\textbf{Condition} & \textbf{Storage} & \textbf{Time points (months)} & \textbf{Attributes monitored / purpose} \\
\midrule
Long-term (ICH Q1A) & 25 C / 60 percent RH & 0, 3, 6, 9, 12, 18, 24 & Appearance, assay, degradants, dissolution, water; supports assigned shelf life \\
Intermediate (ICH Q1A) & 30 C / 65 percent RH & 0, 6, 12 & Same panel; triggered if a significant change occurs at accelerated \\
Accelerated (ICH Q1A) & 40 C / 75 percent RH & 0, 1, 2, 3, 6 & Same panel; supports provisional dating and excursion assessment \\
Photostability (ICH Q1B) & ICH option 1 light exposure & Single challenge & Confirms need for light-protective packaging \\
In-use stability & 25 C / 60 percent RH, bottle opened & Per protocol after first opening & Supports the labeled in-use period for the multi-dose bottle \\
Temperature excursion & Per documented excursion & As incurred & Supports continued-use decisions for monitored excursions \\
\bottomrule
\end{tabularx}
\figcaption{Table 10. Drug-product stability program. Stability-indicating
methods are used at every station; the provisional shelf life is supported by the
long-term and accelerated data and is extended by information amendment as
additional long-term data accrue.}
\end{table}

\subsubsection{Placebo Product}

This is an open-label, single-arm, first-in-human study; there is no comparator
arm, no masking of treatment assignment, and consequently no placebo product. The
absence of a placebo is consistent with the dose-finding and early-feasibility
objectives of a Phase 1 investigation and with ICH E9 principles for early-phase
trials \cite{ichpaistduni,phase1guidance}. In accordance with the ReGARDD IND
content map, this subsection is retained for completeness and to preserve the
table-of-contents numbering, and the placebo product is recorded here as not
applicable. The bias-control measures appropriate to an open-label design,
including the masking of the pathologist who determines R0 resection status and
the central blinded radiographic reader, are described in the study protocol and
in $\S$6.1 of this IND rather than through a placebo. Should a future amendment
introduce a randomized or placebo-controlled component, a corresponding placebo
product description, composition, and matching strategy would be submitted by
information amendment at that time.

\subsubsection{Labeling}

The investigational drug product is labeled in accordance with 21 CFR \S 312.6,
which provides that the immediate package of an investigational drug intended for
human use bear a label with the statement ``Caution: New Drug. Limited by Federal
(or United States) law to investigational use,'' and that the label and labeling
bear no false or misleading statement and make no claim that the drug is safe or
effective for the purposes under investigation \cite{phase1guidance}. Because the
study is open-label, the participant-facing label identifies the product and the
assigned once-daily dose level; no information that would unblind a masked outcome
assessor (for example, the masked pathologist or the central radiographic reader)
appears on any label or in any labeling provided to those assessors. The required
content of the immediate-container label and of the supplied labeling is
summarized in Table 9. The label is affixed during cGMP packaging operations and
is verified at the investigational-product pharmacy against the dispensing record
before any unit is released to a participant. The investigational device and its
single-use accessories carry the investigational-device caution statement under
21 CFR part 812 and are not relabeled or repackaged at the site; their labeling is
addressed in the device content of this submission.

\begin{table}[t]
\centering
\caption{Investigational-product label and labeling content for the daraxonrasib
oral tablet under 21 CFR \S 312.6.}
\label{tab:cmc-labeling}
\begin{tabularx}{\textwidth}{L{4.6cm} Y}
\toprule
\textbf{Label element} & \textbf{Content and basis} \\
\midrule
Investigational caution statement & ``Caution: New Drug. Limited by Federal (or United States) law to investigational use'' (21 CFR \S 312.6) \\
Product identity and strength & Daraxonrasib (RMC-6236) oral tablet; 160 mg, 220 mg, or 300 mg as assigned \\
Protocol and IND reference & Protocol number and IND number; sponsor protocol identifier \\
Lot and batch identification & Manufacturing lot number and a unique container identifier for accountability \\
Storage conditions & ``Store at controlled room temperature; protect from light and moisture; keep in original container'' \\
Expiry or retest date & Date supported by the stability program; not dispensed beyond this date \\
Quantity and directions & Net contents and the once-daily directions for use; ``Swallow whole; do not crush, split, or chew'' \\
Child-resistant statement & ``Keep out of reach of children''; child-resistant closure per special packaging \\
Sponsor identification & Name and address of the sponsor-investigator (ChemicalQDevice) \\
Unblinding control & No information appearing on labeling provided to masked assessors that would reveal dose level or outcome \\
\bottomrule
\end{tabularx}
\figcaption{Table 9. Label and labeling content. The immediate package bears the
21 CFR \S 312.6 investigational caution statement and makes no claim of safety or
effectiveness; the label supports drug accountability and protects the masked
outcome assessors of this open-label study.}
\end{table}

\subsection{Environmental Assessment}

The sponsor claims a categorical exclusion from the requirement to prepare an
environmental assessment or an environmental impact statement under 21 CFR \S
25.31(e), which categorically excludes from environmental review the action of
issuing an investigational exemption for a new drug intended to be administered to
human subjects, provided that the conditions of 21 CFR \S 25.15(c) and \S
25.15(d) are met and provided that no extraordinary circumstances exist that would
result in a significant environmental effect. The investigation described in this
IND meets these conditions. Daraxonrasib will be administered orally to a small
investigational population of up to 18 treated participants at a single academic
site, at clinical doses of 160 to 300 mg once daily over a limited perioperative
schedule; the quantities of drug substance and drug product manufactured,
distributed, and used are correspondingly small and are confined to the
controlled clinical setting and its regulated pharmaceutical-waste stream. The
estimated concentration of the drug substance entering the environment at the
expected use level is below the threshold at which an environmental assessment is
ordinarily required, and no extraordinary circumstances are present: the action
does not significantly affect the quality of the human environment, does not
involve a substance that is persistent, bioaccumulative, or otherwise of
recognized environmental concern at the quantities used, and does not affect any
species or habitat protected under federal environmental law
\cite{phase1guidance}. Investigational product that is unused, returned,
expired, or otherwise unsuitable is destroyed only after sponsor authorization
and in accordance with applicable federal, state, and local regulations governing
the disposal of pharmaceutical waste, and the destruction is documented in the
drug-accountability record. To the sponsor's knowledge, no data or information
exist that would indicate that the proposed investigation may significantly
affect the quality of the human environment. Accordingly, neither an
environmental assessment nor an environmental impact statement is required for
this submission, and the categorical exclusion under 21 CFR \S 25.31(e) is
claimed. Should the scope of the investigation expand in a manner that would
exceed the conditions of the categorical exclusion, an environmental assessment
under 21 CFR \S 25.40 would be submitted by information amendment before
implementing the change.
